\section{ESCL$^*$ Algorithm}\label{sec:SCLAlg}
	In this section, we introduce the ESCL$^*$ algorithm as an extension of SCL$^*$\cite{samadi}. The  ESCL$^*$  algorithm compositionally learns an unknown system $M = \parallel_{i=1}^n M_i$, which consists of $n$ parallel synchronous components. 
	This approach identifies input alphabets and synchronizing actions on-the-fly.  It is initially assumed that each component has an input alphabet 
	and 
	the set of synchronizing actions  $\it sync$ is equal to the union of the input alphabets of all components. Unlike the original SCL$^*$ algorithm,  ESCL$^*$ allows synchronizing actions to have different outputs across the system. However, these outputs must be identical at the point of communication. This assumption avoids
	output non-determinism: the synchronous parallel composition of two deterministic Mealy machines remains deterministic.
	

Let $I^F = \{I_1,\dots, I_n\}$  denote a cover of the total system’s input alphabet where each $I_i\in I^F$ corresponds to a component $i$. If two sets 
$I_i$ and $I_j$ 
 have a non-empty intersection, the corresponding components synchronize on the shared actions. Our goal is to construct such a cover $I^F$.


The algorithm begins with an initial cover consisting of singleton sets, each containing a single action. Suppose the teacher returns a counter-example 
$u_1\ldots u_m$
 in response to an equivalence query. By analysing this counter-example, %it is initially assumed that all actions within the counter-example belong to a single input set. 
 the algorithm assumes all actions within the counter-example belong to a single input set, denoted as $I_{\mathit{CE}}$.
 The algorithm then identifies existing sets in $I^F$ that intersect with $I_{\mathit{CE}}$ to classify shared actions as potential synchronizations.
If an action is correctly identified as a synchronization action, the synchronous composition of the corresponding components must be valid.
An invalid composition indicates that the synchronizing actions were misidentified. This occurs specifically when a synchronizing action has multiple distinct outputs at a synchronization point.
In such cases, the algorithm merges the input sets containing those actions and $I^F$ is updated.
 After updating 
$I^F$, the algorithm relearns the system using  projected alphabet for each 
$I_k\in I^F$. This iterative refinement of the synchronous parallel composition continues until the composed system is equivalent to the original and no further counter-examples are returned.

%We will prove (Proposition~\ref{pro:noallsync}) that \rv{the counter-example $u_1\ldots u_m$ consists of at least one non-synchronizing action} %it is not possible that both actions $u_\ell$ and $u_r$ be the synchronizing actions 
%as the Mealy machine is deterministic and the outputs of synchronizing actions are also unique. 
%Therefore, a sequence of synchronizing actions in the SUL and the learned model will always generate the same sequence of outputs.

 %In the following, we first explain the proposed algorithm by an example and then explain the algorithm in detail.
% If the teacher returns a counter-example, we update $I^F$ and  $\it sync$ as explained.

%\begin{figure}[htb]
%		\centering

 % \subfloat[\label{a}]{\scalebox{0.5}{\includegraphics{example 2.tex}}} 
  %\qquad
  %\subfloat[\label{b}]{\scalebox{0.65}{\includegraphics{ProjectionExample2.tex}}} 
  %\qquad
  %\subfloat[\label{b}]{\scalebox{0.65}{\includegraphics{project2example2.tex}}} 
  
  %\caption{(a) Initial system model $M$ with the input alphabet $\Sigma=\{a,b,c\}$, (b) The projection of $M$ on $\Sigma$ and learn each component individually, (c)  Use the counter-example $ab$ to merge two components}
	%	\label{fig:motivatedexam2}
	%\end{figure}



%	\begin{example}
		%Consider a system model in \figurename{~\ref{fig:motivatedexam2}}(a) whose alphabet is $\{a,b,c\}$. Initially, assume that $a$ is the synchronizing action, i.e., ${\it sync}=\{a\}$, and
  %$I^F=\{\{a\},\{b\},\{c\}\}$. The %teacher may return a counter-example ${\it CE}=ab$ as this string generates the output $11$ in (a) but the output sequence in (b) is $10$. Here, the singelton sets $\{a\}$ and $\{b\}$ are dependent on each other and also $\{a\}\subseteq {\it sync}$. Hence, we merge these singleton sets and update $I^F$ which results in $I^F=\{\{a,b\},\{c\}\}$. As the learning process continues,  the teacher may return a counter-example ${\it CE}=ac$ which generates the output $11$ in (a) but the output sequence is $10$ in (c). Here, the input sets $\{a,b\}$ and $\{c\}$ depend on each other regarding the synchronizing action $a$. However, we do not merge $\{a,b\}$ with $\{c\}$ as they may belong to two distinct components synchronizing on $a$. Instead, we merge  $\{a\}\subseteq {\it sync}$ with $\{c\}\in I^F$
	%	which results in $I^F=\{\{a,c\},\{a,b\}\}$. Now, if the teacher returns another counter-example ${\it CE}=cb$, we find that $\{a,b\},\{b,c\}$  must be merged as they depend on each other regarding the non-synchronizing  actions $b$ and $c$. In this case, we update the input sets, i.e., $I^F=\{\{a,b,c\}\}$, which denotes one component with the alphabet $\{a,b,c\}$.
	%\end{example}
In the following, we provide a detailed description of the ESCL$^*$ algorithm for learning systems composed of synchronous components. In this approach, each component is modelled as a simple Mealy machine. To begin, we first present the formal definitions required for the algorithm's description.


	\begin{definition}[LearnSyncInParts]\label{def:learnparts}
	The \textit{LearnSyncInParts} function takes a Mealy machine $M = (S, s_0, I,O, \delta, \lambda)$ and an input alphabet
	cover $I^F =\{I_1,\dots, I_n\}$, and returns the synchronous parallel composition of learned components.
	% where the synchronous set is defined as ${\it sync}=\bigcup_{ i,j\, |I_i|>1,|I_j|>1,i\neq j} I_i\cap I_j$.
	\[{\it LearnSyncInParts}(M,I^F)=L^*(M,I_1)\parallel\dots\parallel L^*(M,I_n)\]
	%where ${\it sync}=\bigcup_{\forall I_i,I_j\in I^F} I_i\cap I_j$
\end{definition}


%\begin{definition}[SyncActions]\label{def:syncalphabet}
%	The  $\it SyncActions$ function gets a counter-example CE  and a synchronous  set $\it sync$, and returns a set of singleton sets including synchronous actions that are involved in CE:
%	\[{\it SyncActions}({\it CE},{\it sync} )=\{I_j\mid I_j\subseteq {\it sync}, |I_j|=1, ,\exists i\cdot{\it CE}[i]\in I_j\}\]  
%	where $|I_j|$ denotes the size of $I_j$ and the $i^{th}$  action of CE is denoted by ${\it CE}[i]$.
%\end{definition}
%
%\begin{example}
%	If ${\it sync}=\{b,d\}$, and ${\it CE}=bc$, then ${\it SyncActions}({\it CE},{\it sync})=\{\{b\}\}$.
%\end{example}

\begin{definition}[SymbolsSet] \label{def:symbolset}
	The function  $\it SymbolsSet$  takes a counter-example CE and returns a set containing  all symbols present  in CE:
	
	
    	\[{\it SymbolsSet}({\it CE})=\{
	\bigcup_{i=1}^{|{\it CE}|}{{\it CE}[i]}\} \]
  %\rv{, where ${\it CE}[i]$ refers to the $i^{th}$ symbol of CE.}
  , where $|{\it CE}|$ denotes the length of the counter-example.
\end{definition}

\begin{example*}
		If ${\it CE}=abc$, then  ${\it SymbolsSet}({\it CE})=\{a,b,c\}$.
	\end{example*}
\begin{definition}[DependentSets] \label{def:dependset}
	The function  $\it DependentSets$  takes an input cover
	$I^F$, and a set of actions $\it s$, and returns the input sets from $I^F$ that contain at least one  action of the set $s$:	
	\[{\it DependentSets}(I^F,s)=
	\{I_j\mid I_j\in I^F,\exists a\in s\wedge a\in I_j \}\]
 
 %\rv{, where ${\it CE}[i]$ refers to the $i^{th}$ symbol of CE.}
\end{definition}

%	\begin{definition}[InvolvedSets(\textbf{Another!})]
	%	\[{\it InvolvedSets}({\it CE},I^F,{\it sync}^F)=\]
	%	\[\{I_j\mid I_j\in {\it sync}^F ,\exists i\cdot{\it CE}[i]\in I_j\}\,\cup\]  
	%	\[\{I_j\mid I_j\in I^F,\exists i\cdot{\it CE}[i]\not\in {\it sync}\wedge {\it CE}[i]\in I_j \}\]
	%	where the $i^{th}$  character of CE is denoted by ${\it CE}[i]$, and if ${\it sync}=\{s_1,\ldots,s_n\}$ then ${\it sync}^F=\{\{s_1\},\ldots,\{s_n\}\}$.
	%\end{definition}
	
	\begin{example*}
		If $I^F=\{\{a,b\},\{c,d\},\{e,f\}\}$,  $s=\{b,d\}$, then  ${\it DependentSets}(I^F,s)=\{\{a,b\},\{c,d\}\}$.
	\end{example*}
	\begin{definition}[Merge] 	The function ${\it Merge}$	takes two covers $I^{F}_1$ and $I^{F}_2$, and updates $I^{F}_1$ according to the sets in $I^{F}_2$. To this end, it removes the sets of $I^{F}_2$ from $I^{F}_1$ and adds their union to $I^{F}_1$:
		  ${\it Merge}(I^{F}_1,I^{F}_2)= 
		(I^{F}_1\setminus\  I^{F}_2)\cup\{\bigcup_{I_j\in I^{F}_2} I_j\}$
	\end{definition}
	\begin{example*}
		If $I^{F}_1=\{\{a\},\{b\},\{c,d\}\}$, $I^{F}_2=\{\{b\},\{c,d\}\}$ then ${\it Merge}(I^{F}_1,I^{F}_2)=\{\{a\},\{b,c,d\}\}$
	\end{example*}
	
%	\begin{definition}[UnionSyncSets] The $\it UnionSyncSets$ function gets a  cover $I^F = \{I_1, \dots , I_n\}$ and a synchronous alphabet `$m$', and returns the union of the input sets that have a common alphabet as `$m$'.
%		\[{\it UnionSyncSets}(I^F,m)=
%		\{\bigcup_{i:1\dots j} I_i\mid I_i\in I^F,\exists\, j>1\cdot \bigcap_{i:1\dots j} I_i=\{m\}\}
%		\]
%	\end{definition}
%	
%	\begin{example}
%		If $I^F=\{\{a,b\},\{c,b\},\{b\}\}$, then ${\it UnionSyncSets}(I^F,b)=\{a,b,c\}$
%	\end{example}
%


	\begin{algorithm}[tb]
		\small
		\caption{Extended Synchronous Compositional Learning Algorithm}
		\label{alg:SCLalg}
		\hspace*{\algorithmicindent}
		\textbf{Input:} $I^F=\{I_1,\dots,I_n\},M$\\
		\hspace*{\algorithmicindent}
		\textbf{Result:}  $\mathscr{H}$
		\begin{algorithmic}[1]
		\State ${\it sync}:=\bigcup_ {i:1\dots n} I_i \, {\it where}\, I_i\in I^F$ \label{alg:sync}
  		\State $\mathscr{H}\leftarrow {\it LearnSyncInParts}(M,I^F)$\label{alg:Lerarnsync}
			\State ${\it CE}\leftarrow {\it EquivalenceQuery} (\mathscr{H},M)$\label{alg:CEequi}
			\While{${\it CE}\neq {\it yes}$}\label{alg:while}
			\State ${\it ProcessCE}({\it CE})$\label{alg:processCE}
			\State $\mathscr{H}\leftarrow {\it LearnSyncInParts}(M,I^F)$\label{alg:secondLearn}
			\State ${\it CE}\leftarrow {\it EquivalenceQuery}(\mathscr{H},M)$\label{alg:CEequisecond}
			\EndWhile
   \State ${\it sync}:={\it sync}\setminus\{u\mid \nexists\, I_i,I_j\in I^F\cdot u\in I_i \wedge u\in I_j\}$\label{alg:updatesyncsecnd}
			\State ${\it return}\,\mathscr{H}$   \label{alg:return}
		\end{algorithmic}
	\end{algorithm}


	Algorithm \ref{alg:SCLalg} shows the pseudo-code of learning a system compositionally including components that synchronized on some actions. The algorithm
begins with an initial cover $I^F$ consisting of singleton sets for each action in the alphabet $I = \{u_1, u_2, \dots, u_n\}$, such that $I^F = \{\{u_1\}, \{u_2\},\dots, \{u_n\}\}$.
The set of synchronizing actions, $\textit{sync}$, is initially defined as the union of all input alphabets in $I^F$ (line \ref{alg:sync}).
The \textit{LearnSyncInParts} method, then learns each component using the $L^*$ algorithm and returns their synchronous parallel composition (line \ref{alg:Lerarnsync}).
Regarding the learned components, the hypothetical synchronizing actions, \textit{sync} is updated. 
The teacher is asked whether the hypothesis $\mathscr{H}$ is equivalent to $M$ (line \ref{alg:CEequi}).
	If the teacher returns ``yes'' for this equivalence query,
 the algorithm terminates and returns $\mathscr{H}$.
 Otherwise, the counter-example is processed via the \textit{ProcessCE} procedure (line \ref{alg:processCE}) detailed in Algorithm \ref{alg:processCEAlg}, and an(other) iteration of the loop is performed until a correct hypothesis $\mathscr{H}$ is successfully learned (lines \ref{alg:processCE}--\ref{alg:CEequisecond}).
   The $\it sync$ set is refined by removing any actions that are no longer shared among the sets in $I^F$ (line \ref{alg:updatesyncsecnd}).

\begin{algorithm}[tb]
		\small
		\caption{Process Counter-Example}
		\label{alg:processCEAlg}
  \hspace*{\algorithmicindent}
		\textbf{Input:}  ${\it CE}$\\
		\hspace*{\algorithmicindent}
		\textbf{Effect:} Updating the global input cover $I^F$
		\begin{algorithmic}[1]
			\Procedure {$\it ProcessCE$}{}
			\State $I^D\leftarrow \emptyset$
            \State $I_{\it CE}\leftarrow {\it SymbolsSet}({\it CE})$ \label{algline:symbolsset}
            \If {$\exists \, I_i\in I^F\cdot I_i\subseteq I_{\it CE}$} {$I^F\leftarrow I^F\setminus I_i$} \label{algline:updateIF}
            \EndIf
            \State $I^F\leftarrow I^F\cup I_{\it CE}$ \label{algline:addICE}
            \ForEach{$I_j\neq I_{\it CE}\in I^F\cdot I_{\it CE}\cap I_j\neq \emptyset$}~~~~~
            {$I^D\leftarrow I^D\cup \{I_j\}$} \label{algline:generetaeID}
            \EndFor
            \If {$I^D\neq \emptyset$} {$I^D\leftarrow {\it CheckSynSets}(I^D\cup \{I_{\it CE}\})$} \label{algline:learnsynset}   
            \EndIf
            \ForEach{$I_i\in I^D$}~~~ {$I^F\leftarrow {\it Merge}(I^F,I_i)$}\label{algline:merge}
            %\If {$I^D\neq \emptyset$} 
            %\State $I^{D'}\leftarrow{\it IntersectSets}(I^D)$            
             %\State {$I^F\leftarrow {\it Merge}(I^F,\{I_i,I_j\})$}\label{algline:merge}
            \EndFor                     
            \label{alg:secondLearn}
			\EndProcedure
			
		\end{algorithmic}
	\end{algorithm}

 
 		Algorithm \ref{alg:processCEAlg} \rv{is a procedure that} updates the \rv{global variable} %cover 
   $I^F$ using a counter-example CE.
The function \textit{SymbolsSet} returns a set $I_{\it CE}$ containing the actions of the counter-example (line \ref{algline:symbolsset}).
If there exists a set in $I^F$ that is subset of $I_{\it CE}$ (line \ref{algline:updateIF}), then
the algorithm removes it and adds $I_{CE}$ to the cover $I^F$ (line \ref{algline:addICE}). The algorithm then identifies all sets in $I^F$ that share actions with $I_{CE}$ (line \ref{algline:generetaeID}). It is critical to verify whether these shared actions are valid synchronizations actions using Algorithm \ref{alg:learnsync} (line \ref{algline:learnsynset}). Algorithm \ref{alg:learnsync} identifies actions that were incorrectly classified as synchronizing.
 According to the Proposition 1, 
 if these shared actions do not truly belong  to the synchronizing actions set, the corresponding input sets are then merged (line \ref{algline:merge}).



	Algorithm \ref{alg:learnsync} takes a set of alphabet sets containing shared actions and checks whether these  actions are true synchronizing actions. The \emph{LeanSyncInParts} method (line \ref{algline3:learnsyncinpart}), learns each component via the L* algorithm and returns their synchronous composition. If shared actions have conflicting outputs at  the synchronization points, it is concluded that they are not true synchronizing actions
     The \emph{CheckOutput} method (line \ref{algline3:checkoutput}) inspects the synchronization points for output consistency.
     If a synchronizing action has multiple distinct outputs, the algorithm identifies it as an incorrectly identified synchronizing action.
   For each of such invalid shared actions,  the \emph{DependentSets} method finds the input sets containing that action (lines \ref{algline3:forwrongsync}--\ref{algline3:updateID}). Finally, the algorithm returns all input sets that share these invalid synchronizing actions (line \ref{algline3:return}). 
     
 

\begin{algorithm}[thb]
		\small
		\caption{Check Input Sets with Shared Actions}
		\label{alg:learnsync}
        \hspace*{\algorithmicindent}
		\textbf{Input:}  $I^D$:   A set of input sets with shared actions  \\
 %\textcolor{gray}{\text{~~~~//\,\,updated independent covers}}$\\
  \hspace*{\algorithmicindent}
%\textcolor{gray}}$%{\text{~~~~//\,\,independent covers}}$
%		\hspace*{\algorithmicindent}
%		\textbf{Effect:} Updating the set of $I^D$
        \textbf{Result:}  Updated $I^D$ 
		\begin{algorithmic}[1]
			\Function {$\it CheckSyncSets$}{}
			\State $I^M\leftarrow \emptyset,\, I^{D'}\leftarrow\emptyset$ 
            \State $\mathscr{H}\leftarrow {
            \it LearnSyncInParts} (M,I^D)$ \label{algline3:learnsyncinpart}
            \State ${\it wrongsync}\leftarrow {\it CheckOutputs}(\mathscr{H})$ \label{algline3:checkoutput}
            \If {${\it wrongsync}\neq \emptyset$} \label{algline3:Ifwrongsync}
            \ForEach{$u\in {\it wrongsync}$} \label{algline3:forwrongsync}
            \State $I^{M}\leftarrow {\it DependentSets}(I^D,{\it wrongsync})$ \label{algline3:dependentset}
            \State $I^{D'}\leftarrow I^{D'} \cup \{I^M\}$ \label{algline3:updateID}
            \EndFor
            \EndIf
            \State ${\it return}\, I^{D'}$
            \label{algline3:return}
            \label{alg:leransyncset}
			\EndFunction			
		\end{algorithmic}
	\end{algorithm}
 

	
	
	



\begin{proposition}\label{pro:setmerged}
	Let $I^F$ be an input cover, and $\it sync$ be the set of synchronizing actions. If two input sets $I_i,I_j\in I^F$ have a shared action $u$, and $u$ is identified as a non-synchronizing action  ($u \notin sync$), then %\rv{$I^F$ must be updated according to    the union of $I_i$ and $I_j$.} 
 $I_i$ and $I_j$ must be merged.
\end{proposition}

\begin{proof}
	
		Assume that in Algorithm~\ref{alg:processCEAlg}, %lines \ref{alg2:If}-\ref{alg2:updateOutSync}, 
        $I^F$ contains two input sets  $I_i=\{a_1,\ldots,a_n,u\}$ and $I_j=\{b_1,\ldots, b_m,u\}$, where $u$ is detected as a non-synchronizing action. By Definition~\ref{def:dependactions}, 
	each input set consists of dependent actions. Hence, the actions $a_1,\ldots,a_n$ are dependent on the non-synchronizing action $u$, and $b_1,\ldots,b_m$  are likewise dependent on $u$. Since $u$ is a non-synchronizing action, the transitive property of the dependency relationship implies that $a_1,\ldots,a_n$ are also dependent on $b_1,\ldots,b_m$. %, and so,  $I_i$ and $I_j$ must be merged.
 \rv{As a result, $I_i$ and $I_j$ are removed from $I^F$ and instead their union $\{a_1,\ldots,a_n,b_1,\ldots,b_m,u\}$ is added to $I^F$.}
\end{proof}

\subsection{Minimum Counter-Examples}
\rv{We aim to find a counter-example with minimal length where any of its prefixes are not a counter-example.}
If the teacher returns a non-minimal counter-example, some input sets may be merged incorrectly~\cite{30}. 
To address this, we extend the algorithm proposed in \cite{30} to minimize the returned counter-example based on the synchronizing actions. Our approach iteratively identifies the smallest prefix of $CE$ that remains a valid counter-example.   To this end, we select the input sets  $I^M$ including synchronizing and non-synchronizing actions participating in  CE. Then, we iteratively select a subset of $I^M$, merge its members, and produce a set $A$. The projection of CE on $A$ is returned as a minimum counter-example if the SUL and the hypothesis model produce
different outputs.% for it.
%{\it MinimizeCE} (Algorithm~\ref{alg:CEMinimize}) we iteratively get a subset of ${\it DependentSets}({\it CE}, {I^F},{\it sync})$ and merge its actions. % which leads to producing a set $A$.
%The algorithm returns a projection of a counter-example %, i.e., $P_{A}(\it CE)$,  as an output if it is a counter-example.
%In Algorithm~\ref{alg:CEMinimize}, line \ref{alg3:Minprefix},  the function \emph{MinPrefixCE} takes a counter-example CE and returns the smallest prefix of CE that is also a counter-example. On lines \ref{alg3:updateI^D}-\ref{alg3:foerach1}, the input covers including synchronizing and non-synchronizing actions participating in the counter-example CE are selected. On lines \ref{alg3:foreach2}-\ref{alg3:A}, we iteratively get a subset of $I^D$, merge its members, and produce a set $A$. On lines \ref{alg3:CEmin}-\ref{alg3:If},  the new counter-example is returned as a  minimum counter-example if the SUL and the hypothesis model produce different outputs for it.
%\begin{algorithm}[tb]
%		\small
%		\caption{Counter-Example Minimization}
%		\label{alg:CEMinimize}
%  \hspace*{\algorithmicindent}
%  \textbf{Input:} ${\it CE},I^F=\{I_1,\dots,I_n\},M,\mathscr{H},{\it sync}$\\
 % \textbf{Result:} ${\it CE}_{\it Min}$	
%		\begin{algorithmic}[1]
%			\Procedure {$\it MinimizeCE$}{}
%			\State ${\it CE}\leftarrow {\it MinPrefixCE}({\it CE})$\label{alg3:Minprefix}
 %           \State $I^D\leftarrow{\it DependentSet}({\it CE},I^F,{\it sync})$\label{alg3:updateI^D}
  %          \ForEach{$s\in {\it sync}\,\textbf{s.t}\, \exists i\cdot{\it CE}[i]=s$ }   $I^D\leftarrow I^D \cup\{\{s\}\}$\label{alg3:foerach1}
   %         \ForEach{$k\in\{2,\ldots,|I^D|\}$}\label{alg3:foreach2}
    %        \State $I^*\leftarrow\{I_j\mid |I_j|=k\,\wedge\,\forall\,u\in I_j\rightarrow u\in I^D\}$\label{alg3:I*}
            %I_j=\{I_1,\ldots,I_k\}\,\textbf{s.t}\, \forall\,i\in\{1,\ldots,k\}\cdot\, I_i\in I^D\}$
     %       \ForEach{$I_\rho\in I^*$}\label{alg3:foreach3}
      %      \State $A\leftarrow\bigcup_{v\in I_\rho} I_\rho$\label{alg3:A}
       %     \State ${\it CE}_{\it Min}\leftarrow P_A({\it CE})$\label{alg3:CEmin}
        %    \If{${\it CE}_{\it Min} \textit{is a counter-example}$}  ${\it return} {\it CE}_{\it Min}$\label{alg3:If}
         %   \EndIf
          %  \EndFor
           % \EndFor
            %\EndFor
			%\EndProcedure
		%\end{algorithmic}
	%\end{algorithm}
	In the next section, we prove that the ESCL$^*$ algorithm terminates and the synchronous components are learned correctly.
